\begin{frame}{6.1 정리}
  \begin{itemize}
    \item 정규화는 손실함수에 ``가중치 크기'' 페널티 한 항을 더해 유효 유연성을
      낮춤 — 4.1의 언어로 $\lambda$를 올리면 \textbf{분산을 편향과 바꾸는} 것.
      살짝만 올려도 Var가 4자릿수 급감할 만큼 ``분산 제거''에 효율적.
    \item L2는 모든 가중치를 \textbf{매끄럽게} 축소, L1은
      soft-thresholding($|z_j|\le\lambda$이면 0)으로 약한 특징을
      \textbf{정확히 0}에 스냅시켜 특징 선택을 \emph{무상}으로 —
      기하학적으로 ``원 위'' vs ``마름모 꼭짓점 위''.
    \item bias는 관례적으로 정규화 대상에서 \textbf{제외}.
  \end{itemize}
  \vskip0.4em
  \kq{``$\lambda$ 자체는 얼마가 적정한가'' — U자형 곡선의 바닥을 정직하게
  고르는 절차는 6.2(교차검증)·6.3(train/val/test)에서.}
\end{frame}
